\documentclass[11pt]{article}
\usepackage{amsmath,amssymb,amsthm,hyperref}
\newtheorem{prop}{Proposition}
\newcommand{\sets}[1]{\{#1\}}
\title{A timely clearance certificate for partner substitution}
\date{24 November 2026}
\begin{document}
\maketitle

An instrument reading all $n$ service lanes can be replaced by one reading only $r+1$, where $r$ prior operations still hold distinct lanes. The installed receiver and its safety and deadline proof survive the change even when the chosen lane changes. A third instrument produces the same certificate as the replacement, but too late. It passes a standalone deadline and fails the joint one. The minimum number of distinct timely reports is exactly $r+1$, rather than in general the number of possible old occupancy states.

\section{The shared work and the decision cut}
Lanes are separate unit-capacity resources, named $0,\ldots,n-1$, with $1\le r<n$. At tick 0 the original admission desk accepts one old maintenance job in each of $r$ distinct lanes. Every sequence of $r$ distinct valid lane selections is admitted; any $r$-set $B$ therefore occurs. Each old job acquires a non-preemptible physical reservation on its lane for $[0,5)$ and completes at tick 5. At tick 1 the outgoing instrument loses its right to issue new reports and the new instrument is installed. This change neither revokes nor moves old work. The clearance instrument cannot cancel old jobs, reserve lanes, or change their deadlines.

There is one new service due by tick 4. It takes two whole ticks on one lane and must dispatch by tick 2. A resource rule refuses dispatch on an occupied lane. The receiver submits one irrevocable lane choice at tick 2; it cannot probe occupancy, test an alternative, or delay. The refusal is a safe service failure, not an additional decision channel. No further old work can enter before tick 5, and no other new work competes for a lane. All old jobs complete on time when the safety rule is obeyed. The standalone comparison is the same service with a deadline at tick 5 instead, without the incumbent receiver's earlier joint cut. It measures safe completion and the eventual certificate, not equality of every timestamp.

At tick 2 the receiver knows the common schedule and the capacity promise $|B|=r$, but neither $B$ nor a private sensor read. Apart from the instrument's one fixed-cut report, all decision-time observations are the same on all $B$. Message contents and observable availability are counted as distinct observations; latency, silence or external cues cannot encode an uncounted distinction. The lower bound below applies to this information structure, not to a receiver with its own occupied-lane sensor or the ability to retry before the cut. Both time intervals are half-open.

\section{Interface and replacement}
The installed interface is a trusted clearance assertion $(i,[2,5))$ delivered no later than tick 2. Its meaning is that lane $i$ has no old reservation during that interval. A conventional encoded tag can be decoded to its fixed lane name and interval; the lane is the variable part. The installed receiver accepts the clearance if current and covering its service interval, dispatches on $i$, and otherwise stays. Its existing proof has three inputs:
\[
  \text{a receipt at the cut},\qquad
  i\notin B,\qquad
  \text{no further conflicting reservation in the asserted interval}.
\]
The first input enables the unchanged receiver action. The last two preserve the old reservations and permit the new service to finish at tick 4. Independently, old reservations and their physical releases at tick 5 are invariant under any valid new action. The instrument cannot supply that invariant merely by asserting its own correctness.

The original full-grid instrument reads each lane and asserts the highest clear index. The replacement has only a bank of $r+1$ lane sensors, on fixed indices $0,\ldots,r$. It asserts its lowest clear index. The per-lane observation law is exact at tick 2, and the old reservation law precludes a conflicting arrival or reassignment through tick 5. There must be a clear bank lane, even if every hidden old job is in an inconvenient location. Those two facts derive the clearance meaning and timing without reading the rest of the grid. The receiver's code, proof, due time, resource rule and old-work completion condition do not change. The original and new chosen actions may differ. The proof is uniform over all sound current assertions, not a proof specifically for the original action.

For example at $n=4,r=1$, legacy occupancy $B=\{2\}$ and $B'=\{3\}$ both produce the same bank read and early certificate for lane 0. The replacement does not recover the original observations. When $n\ge r+3$, a bank reading with $r-1$ occupied bank lanes similarly leaves at least two distinct possibilities outside the bank. This compression is lossless for delivering a valid certificate, not for answering unrestricted occupancy queries. Exact per-lane sensing, a synchronized delivery phase and the immutable reservation interval are model assumptions. A point-in-time idle reading does not establish clearance by itself.

\section{Exact timely information requirement}
\begin{prop}
For a uniform one-message policy with fixed public decision context, the minimum number of distinguishable timely observations needed to guarantee the joint service and preserve every old reservation is $r+1$. For an actually reachable decision cell, a deterministic causal choice exists exactly when there is a lane outside the union of all its possible old occupancy sets.
\end{prop}
\begin{proof}
The new bank is one construction with $r+1$ observations. The announced lane determines the service action; every one of its possible outputs can occur. Conversely, consider a receiver facing $k\le r$ observations. Each released observation that can meet the deadline selects one lane; an observation without dispatch cannot help. An old set $B$ of size $r$ can include all those at most $k$ selected lanes, and is reachable by the ordinary old admission rule. Every choice then misses service, even if its damaging effect is prevented by the resource rule. For a specified observation cell, its shared action addresses every $B$ in the cell exactly when that lane is outside their union. This is both a necessary action rule and a construction of a uniform decoder. All successful continuations then follow from the unchanged resource law.

Randomization under an almost-sure guarantee does not relax the bound. The finitely many initial $B$ have a common measure-one set of seeds on which the policy succeeds for every $B$ simultaneously. Fixing any seed in that set gives a deterministic encoder and decoder on the same alphabet. The seed distribution must be independent of $B$ except for the modeled observations; partially hidden correlations are not supplied as public information. With a positive tolerated failure probability this argument no longer proves the same bound.
\end{proof}

The delayed instrument samples exactly the replacement's bank and sends exactly its report at tick 3. It can dispatch there and finish at tick 5 under the standalone deadline. Under the installed joint deadline it has one uninformative decision view. Every complete legacy history has a correct choice in hindsight, and that very choice arrives at tick 3. There is no function of the information available at tick 2 that chooses it correctly across all histories: for each guessed lane there is an initial $B$ that holds it. Waiting is safe but misses the service. The fixed incumbent acts only on a current certificate and does not dispatch. Thus stand-alone completion and even equality of eventual messages cannot justify this substitution.

For $r=1$, any proper subfamily of the $n$ no-report histories admits a common action, but the whole family does not. Even checking all $(n-1)$-wise combinations therefore cannot replace checking the full information cell. The obstruction is collective; a pairwise conflict graph does not detect it for $n\ge3$. This claim concerns intersections of viable actions, rather than reconstructing a unique required action.

\section{When the service needs two lanes}
If one admission needs an atomic bundle of $b$ distinct lanes for the same interval, each distinct decoded message chooses a $b$-set $A$. Assuming $n\ge r+b$, an exact condition for a menu $\mathcal A$ is
\[
  \forall B\in\binom{[n]}r\ \exists A\in\mathcal A:\quad A\cap B=\varnothing .
\]
Necessity follows because a decoder has no other private evidence; sufficiency follows by choosing on the instrument side any disjoint bundle and deriving clearance on each lane. This characterizes minimum alphabet size, not the sensing cost of computing a valid message. Complements of the menu are an ordinary covering design $C(n,n-b,r)$. A set hitting every $A$ needs more than $r$ points, or a reachable blocking set would contain a hitting set. Thus $|\mathcal A|\ge r+1$, attained by $r+1$ disjoint bundles when $n\ge b(r+1)$. At $n=r+b$ each $B$ has only one complementary bundle, so $\binom nr$ messages are needed.

For $b=2$, view the decoded pairs as edges of a graph on the lanes, including isolated vertices. An $r$-set meets every edge exactly when its $n-r$-element complement is independent. Thus the complement graph must have clique number at most $q=n-r-1$. By Tur\'an's theorem it has at most the edges of the complete $q$-partite graph with balanced part sizes. The optimum menu is its complement: disjoint cliques with balanced sizes $s_1,\ldots,s_q$, of cost $\sum_i\binom{s_i}2$. For $n=5,r=2,b=2$, this gives four reports, whereas a single-lane service needs three. The extra query changes the information requirement even at the same cut. These are standard graph extremality and covering arguments, not new theorems in those fields.

\section{Reuse and direct comparison}
The measurement and delivery checks must be proved for each instrument. There is one unchanged incumbent dispatch argument and one unchanged reservation/noninterference argument. The original local observation proof uses $n$ exact readings and existence of any clear lane; the replacement local proof uses $r+1$ exact bank readings and existence of a clear bank lane. The lower sensor count is a count of inspected lanes in this model, not a runtime or cost claim for a physical instrument. The delayed replacement fails the delivery condition despite satisfying the later value and interval conditions. For either successful instrument, all raw receipt data not consumed by the installed rule or needed to establish its inputs stay outside the receiver's proof. No proof-work advantage over an equally factored direct assume--guarantee argument was demonstrated. That comparator reuses the same two arguments and owes the same sensor, timing and interval checks.

The pinned \href{https://github.com/eternal-november/eos-research-programmes/blob/e1348964f73704212d2d8b48d5982a5ed600a2ff/projects/contextual-composition/feedback-and-concurrency/observable-covers/Residual_Profiles_and_Admission_Reflection.tex}{Residual Profiles and Admission Reflection} factorization is applicable to the installed receiver's primitive test of the delivered lane and window, with the not-yet-arrived observation explicit. The old reservation proof is a separate plant invariant. Exact equality of all raw observations, equality of chosen actions, and fullness for arbitrary new consumers are not transport claims here. There is no native witness comparison as required by that source's stronger common-profile theorem. Each plant history is generated first; its reachable decision cell has to carry a sound certificate for every member, not just some full-trace match under lookahead.

\href{https://web.cs.wpi.edu/~heineman/html/teaching_/CS562/p109-de_alfaro.pdf}{De Alfaro and Henzinger, Interface Automata}, in particular the compatible-input condition for refinement, is the direct temporal interface comparator. \href{https://arxiv.org/abs/2205.12085}{Finkbeiner, Metzger and Moses} already treat distributed information-flow obligations across traces. \href{https://arxiv.org/abs/math/9502238}{Gordon, Kuperberg and Patashnik} use the covering design formulation for fixed-size covering sets. The two-lane formula invokes Tur\'an's theorem (see \href{https://www.renyi.hu/~p_erdos/1971-01.pdf}{Erd\H{o}s's account and its reference to the original result}). There is no priority claim.

The deterministic replay generates 494 reachable cuts across 28 parameter suites, checks 2,964 paired joint and standalone resource traces by replaying raw reservations and service actions, and separately enumerates action menus in small instances. It checks explicit history-specific hindsight actions which cannot be implemented from the common late view. This is a synthetic, fixed-deadline C4 result; it is not acceptance of C4 or C11, a claim about changing the continuation law, clock drift or recourse, or a physical or Lean validation.
\end{document}
